% ================================================================
%  TAREA 1 — Sistema Solar
%  Fecha de asignación: martes, 5 de mayo de 2026
% ================================================================
\chapter{Tarea 1 — Sistema Solar}

\begin{destacado}
    \textbf{Fecha de asignación:} Martes, 5 de mayo de 2026 \quad
    \textbf{Asignada por:} Daniel Soto Villada
\end{destacado}

\medskip

Esta primera tarea está organizada en dos secciones. La primera invita a reflexionar
sobre la escala de tiempo del Sistema Solar en relación con la historia humana; la
segunda profundiza en las distancias y tamaños característicos de los principales
componentes del sistema. Antes de comenzar, asegúrate de leer con atención el
enunciado de cada punto.

\begin{nota}{Contexto previo necesario}{nota:contexto-tarea1}
    Estos problemas se desarrollarán en los encuentros del semillero una vez que se
    haya presentado el marco teórico correspondiente: posición del Sistema Solar en
    la galaxia, estructura interna del sistema (planetas, cinturón de asteroides,
    heliosfera, nube de Oort) y manejo de unidades astronómicas. Se incluyen aquí
    para que los estudiantes puedan revisar los enunciados con anticipación.
\end{nota}

% ---------------------------------------------------------------
\section{El Sistema Solar y la Historia Humana}
% ---------------------------------------------------------------

Desde la aparición de los \textit{Homo sapiens} hace aproximadamente 300\,000 años,
el Sistema Solar ha continuado su órbita alrededor de la Vía Láctea.

Sabiendo que:
\begin{itemize}[leftmargin=1.5em]
    \item el Sistema Solar tarda \textbf{230 millones de años} en dar una vuelta
          completa a la galaxia;
    \item su velocidad orbital es de aproximadamente \textbf{828\,000\,km/h};
\end{itemize}

\begin{actividad}{Preguntas sobre la órbita galáctica}{act:orbita-galactica}
    \begin{enumerate}[leftmargin=1.5em, label=\arabic*.]
        \item ¿Qué fracción de la órbita galáctica ha recorrido el Sistema Solar
              desde la aparición de los humanos modernos?
        \item ¿A qué distancia en kilómetros equivale ese recorrido?
        \item Si la órbita fuera un círculo, ¿qué ángulo (en grados) ha barrido
              el Sistema Solar en ese tiempo?
    \end{enumerate}
\end{actividad}

% ---------------------------------------------------------------
\section{Tamaño y Distancias en el Sistema Solar}
% ---------------------------------------------------------------

\begin{concepto}{Unidades de referencia}{con:unidades-ss}
    \begin{itemize}[leftmargin=1.5em]
        \item $1\,\text{AU} \approx 150 \times 10^{6}\,\text{km}$
        \item \textbf{Heliosfera:} $\in [80,\,100]\,\text{AU}$
        \item \textbf{Nube de Oort:} $\in [5\,000,\,100\,000]\,\text{AU}$
    \end{itemize}
\end{concepto}

\begin{actividad}{Escala del Sistema Solar}{act:escala-ss}
    \begin{enumerate}[leftmargin=1.5em, label=\arabic*.]

        \item Si la órbita de la Tierra tuviera $1\,\text{cm}$ de radio, ¿a qué
              distancia (en centímetros o metros) estaría\ldots
              \begin{enumerate}[leftmargin=1.5em, label=\alph*)]
                  \item la órbita de Neptuno?
                  \item el borde de la Heliosfera?
                  \item el inicio de la Nube de Oort?
                  \item el fin de la Nube de Oort?
              \end{enumerate}

        \item Si la Tierra está a $1\,\text{AU}$ del Sol, ¿cuántas veces más lejos
              está el borde externo de la Nube de Oort comparado con la Tierra?

        \item Si un avión vuela a $900\,\text{km/h}$ de manera constante partiendo
              desde la Tierra en línea recta, ¿cuánto tiempo tardaría en llegar\ldots
              \begin{enumerate}[leftmargin=1.5em, label=\alph*)]
                  \item a la Luna?
                  \item a Neptuno?
                  \item al borde de la Heliosfera?
                  \item a la parte externa de la Nube de Oort?
              \end{enumerate}

        \item Sabemos que un cometa interestelar de órbita parabólica que acaba de
              pasar por el Sol viaja en promedio a $1\,\text{km/s}$.
              \begin{enumerate}[leftmargin=1.5em, label=\alph*)]
                  \item ¿Hace cuánto tiempo entró al Sistema Solar?
                  \item ¿Hace cuánto tiempo pasó por el Cinturón de Kuiper?
                  \item ¿Hace cuánto tiempo pasó por el Cinturón de Asteroides?
                  \item ¿Hace cuánto tiempo pasó por el planeta Tierra?
              \end{enumerate}

        \item ¿Cuántas veces más grande es el volumen de la Nube de Oort en
              comparación con el volumen de la región del Cinturón de Kuiper?

              \begin{destacado}
                  \textbf{Volumen de una esfera de radio $r$:}
                  \[
                      V = \frac{4\pi}{3}r^3
                  \]
              \end{destacado}

        \item La nave \textit{Voyager~1} viaja a $60\,000\,\text{km/h}$.
              ¿Cuántos años tardó en atravesar la Heliosfera? ¿Es este resultado
              consistente con la realidad?

        \item Si la luz tarda 8 minutos en llegar del Sol a la Tierra, ¿cuál es
              la velocidad de la luz? ¿Cuánto tardaría la luz en llegar al borde
              exterior de la Nube de Oort?

    \end{enumerate}
\end{actividad}
